\pdfoutput=1
\documentclass[12pt,a4paper]{article}
\usepackage[english]{babel}
\usepackage[T1]{fontenc}
\usepackage[a4paper,top=2cm,bottom=2cm,left=2.5cm,right=2.5cm]{geometry}
\usepackage{amsmath,amssymb,amsthm}
\usepackage{mathtools}
\usepackage{bm}
\usepackage{bbm}
\usepackage{graphicx}
\usepackage{tikz}
\usepackage{xcolor}
\usepackage{float}
\usepackage{booktabs}
\usepackage{enumitem}
\usepackage{paralist}
\usepackage{algorithm}
\usepackage{algpseudocode}
\usepackage[colorlinks=true, citecolor=blue, urlcolor=blue]{hyperref}
\usepackage{cleveref}

\newtheorem{theorem}{Theorem}[section]
\newtheorem{lemma}[theorem]{Lemma}
\newtheorem{proposition}[theorem]{Proposition}
\newtheorem{corollary}[theorem]{Corollary}
\newtheorem{definition}[theorem]{Definition}
\newtheorem{conjecture}[theorem]{Conjecture}
\newtheorem{assumption}[theorem]{Assumption}
\theoremstyle{remark}
\newtheorem{remark}[theorem]{Remark}
\newtheorem{example}[theorem]{Example}
\newtheorem{protocol}[theorem]{Protocol}

\newcommand{\R}{\mathbb{R}}
\newcommand{\C}{\mathbb{C}}
\newcommand{\N}{\mathbb{N}}
\newcommand{\Z}{\mathbb{Z}}
\newcommand{\Q}{\mathbb{Q}}
\newcommand{\E}{\mathbb{E}}
\newcommand{\Pbb}{\mathbb{P}}
\newcommand{\D}{\mathcal{D}}
\newcommand{\F}{\mathcal{F}}
\newcommand{\G}{\mathcal{G}}
\newcommand{\T}{\mathcal{T}}
\newcommand{\loss}{\mathcal{L}}
\newcommand{\Obs}{\mathcal{O}}
\newcommand{\M}{\mathcal{M}}
\newcommand{\A}{\mathcal{A}}
\newcommand{\B}{\mathcal{B}}
\newcommand{\I}{\mathcal{I}}
\newcommand{\J}{\mathcal{J}}
\newcommand{\X}{\mathcal{X}}
\newcommand{\K}{\mathcal{K}}
\newcommand{\Sg}{\mathcal{S}}
\newcommand{\U}{\mathcal{U}}
\newcommand{\W}{\mathcal{W}}
\newcommand{\AuditGrp}{\mathfrak{G}_{\text{audit}}}
\newcommand{\ExpertSet}{\EuScript{E}}
\newcommand{\CGC}{\mathcal{C}}

\newcommand{\SCX}{\textsf{SCX}}
\newcommand{\Yajie}{\textsf{Yajie}}
\newcommand{\Spring}{\textsf{Spring}}
\newcommand{\Cercis}{\textsf{Cercis}}
\newcommand{\Situs}{\textsf{Situs}}
\newcommand{\MoE}{\textsf{MoE}}
\newcommand{\DeFi}{\textsf{DeFi}}
\newcommand{\MEV}{\textsf{MEV}}

\DeclareMathOperator{\argmax}{arg\,max}
\DeclareMathOperator{\argmin}{arg\,min}
\DeclareMathOperator{\sgn}{sgn}
\DeclareMathOperator{\diag}{diag}
\DeclareMathOperator{\spec}{spec}
\DeclareMathOperator{\softmax}{softmax}
\DeclareMathOperator{\topk}{top\text{-}k}
\DeclareMathOperator{\Cov}{Cov}
\DeclareMathOperator{\Var}{Var}
\DeclareMathOperator{\Tr}{Tr}
\DeclareMathOperator{\TV}{TV}
\DeclareMathOperator{\KL}{KL}
\DeclareMathOperator{\rank}{rank}
\DeclareMathOperator{\TVL}{TVL}

\newcommand{\rigorFull}{\textbf{[Full Proof]}}
\newcommand{\rigorPartial}{\textbf{[Partial Proof]}}
\newcommand{\rigorSketch}{\textbf{[Proof Sketch]}}
\newcommand{\openproblem}{\textbf{[Open Problem]}}
\newcommand{\honestcrit}[1]{\textcolor{red}{\textbf{[Honest Critique]} #1}}

\title{\textbf{Decentralized Finance Audit and Maximal Extractable Value\\
as Gauge Exploitation: An SCX Framework}}
\author{SCX}
\date{\today}

\begin{document}

\maketitle

\begin{abstract}
We present a unified theoretical framework for understanding two foundational
problems in decentralized finance (DeFi) --- maximal extractable value (MEV)
and smart contract security auditing --- as manifestations of the same
underlying phenomenon: \emph{gauge exploitation} in multi-expert decision
systems. Using the State-Conditioned Expertise (SCX) framework originally
developed for label noise detection, we show that MEV extraction is the
identification and exploitation of unaligned coordinate systems (gauges) in
the transaction ordering space of a blockchain. Block proposers exploit the
fact that the mempool ordering is not gauge-fixed; transactions within a block
are gauge-equivalent under reordering, and MEV searchers find gauge
transformations that extract value. Smart contract audits, oracle price feeds,
and flash loan attacks are each re-framed as particular instances of the same
gauge exploitation problem, where the \Cercis{} score serves as a universal
measure of expert (dis)agreement. We prove four theorems bounding MEV profit,
audit detection probability, oracle manipulation detectability, and flash loan
attack windows under the SCX framework. We further diagnose the 2008 financial
crisis credit rating failure as an \emph{under-declared} M=1 expert review
disguised as M=3, offering a structural explanation rooted in gauge
degeneracy rather than mere malfeasance.
\end{abstract}

\section{Introduction}
\label{sec:introduction}

Decentralized finance has grown from a niche experiment into a multi-hundred-billion
dollar ecosystem of lending protocols, decentralized exchanges, stablecoins, and
derivative markets~\cite{schaer2021defi}. Two challenges have persisted throughout
this growth: the extraction of value by block proposers through transaction
reordering (maximal extractable value, or MEV)~\cite{daian2020flashboys}, and
the difficulty of ensuring smart contract security through third-party
audits~\cite{li2020survey}.

At first glance, MEV and smart contract auditing appear unrelated. MEV concerns
the economics of block building in permissionless consensus protocols; auditing
concerns the correctness of code before deployment. Yet both share a deeper
structural similarity: they involve multiple agents making assessments under
partial information, where one agent's advantage comes from exploiting the
difference between local and global perspectives.

In this paper we argue that both problems are instances of \emph{gauge
exploitation}, a concept we formalize within the SCX
framework~\cite{scx2025,yajie2025,spring2025}. The original SCX framework
introduced the \Cercis{} score
\begin{equation}
S(s) = Q(s) + \eta(t) \cdot N(s)
\label{eq:cercis}
\end{equation}
where $Q(s)$ measures the quality (consensus) of a set of expert assessments on
sample $s$, $N(s)$ measures the novelty (disagreement) among those assessments,
and $\eta(t)$ is a time-dependent annealing parameter. The framework was
originally developed for label noise detection in supervised learning, where
$M>1$ expert classifiers vote on the correct label.

We extend this framework to three domains in decentralized finance:
\begin{enumerate}
\item \textbf{MEV as gauge exploitation}: Block builders exploit the
gauge-freedom of transaction ordering within a block. The \Cercis{} score
measures how much value a particular ordering extracts relative to a
fair-ordering baseline.
\item \textbf{Smart contract audits as $M>1$ expert review}: Multiple
independent security firms act as an expert ensemble. Their consensus
increases $Q(s)$; their disagreement signals remaining risk.
\item \textbf{Oracle manipulation as gauge attacks}: Oracles provide external
data; choosing which oracle to trust is a gauge choice. Multi-oracle
consensus detects manipulation through \Cercis{} disagreement.
\end{enumerate}

We further analyze flash loan attacks as temporary gauge misalignments and
diagnose the failure of credit rating agencies during the 2008 financial crisis
as an M=1 under-declared system --- three agencies using effectively the same
model, producing correlated (not independent) assessments.

\subsection{Contributions}

Our contributions are as follows:
\begin{itemize}
\item A formal mapping between MEV extraction and gauge theory in the SCX
framework, with a theorem bounding MEV profit by the \Cercis{} score spread
(Theorem~\ref{thm:mev_bound}).
\item A theorem showing that $M$-auditor smart contract security guarantees scale
as $O(1/\sqrt{M})$ detection probability under independence assumptions
(Theorem~\ref{thm:audit_scaling}).
\item A theorem bounding oracle manipulation detectability through multi-source
\Cercis{} disagreement (Theorem~\ref{thm:oracle_detect}).
\item A theorem bounding the flash loan attack window (Corollary~\ref{cor:flash}).
\item A structural diagnosis of credit rating agency failure as gauge degeneracy,
not fraud.
\item A summary mapping table (Table~\ref{tab:mapping}) connecting all domains
to the SCX framework.
\end{itemize}

\section{Preliminaries}
\label{sec:preliminaries}

\subsection{The SCX Framework}

The State-Conditioned Expertise (SCX) framework~\cite{scx2025} models a
decision system as a collection of $M$ experts, each with a local
\emph{gatekeeper} function $g_i: \X \rightarrow [0,1]$ that maps an input
$s \in \X$ to a confidence score. The \Cercis{} score
(Equation~\ref{eq:cercis}) decomposes into:
\begin{itemize}
\item $Q(s) = \frac{1}{M} \sum_{i=1}^M g_i(s)$, the mean expert confidence
(quality term).
\item $N(s) = \frac{1}{M} \sum_{i=1}^M (g_i(s) - Q(s))^2$, the variance
among experts (novelty term).
\item $\eta(t) \geq 0$, a temperature parameter that decays over time
$t$~\cite{yajie2025}.
\end{itemize}

The key insight is that when $M=1$ (a single expert), $N(s)=0$ identically, and
$S(s)=Q(s)$ --- the system has no notion of uncertainty beyond the single
expert's confidence. This is referred to as an \emph{under-declared} system:
the framework cannot distinguish between high-confidence correctness and
high-confidence error.

\subsection{Gauge Theory Analogy}

In physics, a \emph{gauge theory} is defined by a configuration space where
physically equivalent states are related by a local symmetry transformation
(a \emph{gauge transformation}). The choice of a particular representative
from each equivalence class is called \emph{gauge fixing}. Physical observables
must be gauge-invariant --- they must not depend on which representative is
chosen.

We draw an explicit analogy for DeFi systems:
\begin{itemize}
\item \textbf{Configuration space}: The set of all possible orderings of
pending transactions in the mempool.
\item \textbf{Gauge transformation}: A reordering of transactions that preserves
the block's validity but changes who profits from arbitrage, liquidations, or
front-running.
\item \textbf{Gauge fixing}: The consensus protocol's block proposal rule (e.g.,
highest-gas-price-first, or a fair-ordering protocol).
\item \textbf{MEV}: The value extracted by choosing a gauge (ordering) that
benefits the block proposer at the expense of other users.
\end{itemize}

A \emph{gauge exploitation attack} occurs when an agent finds a gauge
transformation that shifts value from other participants to themselves, and
the protocol lacks an invariant that prevents such transformations.

\subsection{Expert Ensembles in DeFi}

DeFi systems naturally produce expert ensembles in several forms:
\begin{enumerate}
\item \textbf{Audit firms}: Before deploying a smart contract, projects engage
one or more security auditors. Each auditor applies its own methodology,
tooling, and expertise --- effectively, each auditor is an expert with a
different gatekeeper.
\item \textbf{Oracle providers}: Multiple decentralized oracle networks
(Chainlink, Pyth, Tellor, etc.) report the same external data (e.g., an
asset price) on-chain. Each oracle is an expert measuring the same quantity
with different data sources and aggregation methods.
\item \textbf{Block builders}: In proposer-builder separation (PBS), multiple
block builders construct candidate blocks. Each builder orders transactions
differently, and the proposer selects among them. Each builder is an expert
in maximizing their own returns.
\item \textbf{Liquidators}: In lending protocols, multiple liquidators compete
to identify and liquidate under-collateralized positions. Each liquidator
uses a different monitoring strategy and gas bidding approach.
\end{enumerate}

\section{MEV as Gauge Exploitation}
\label{sec:mev}

Maximal extractable value (MEV) refers to the value that a block proposer can
extract by arbitrarily including, excluding, or reordering transactions within
a block~\cite{daian2020flashboys}. Since its identification in Ethereum, MEV
has grown into a multi-billion dollar ecosystem~\cite{qin2022quantifying}.

\subsection{The Mempool as Configuration Space}

Consider the \emph{mempool} $\M_t$ at time $t$: the set of all pending
transactions that have been broadcast but not yet included in a block. A block
$B$ is an ordered subset of $\M_t$ that satisfies the protocol's validity
rules. The set of all valid blocks is the configuration space $\Omega$ of our
gauge theory.

For a given set of transactions $\T \subseteq \M_t$, let $\Pi(\T)$ be the set
of all possible orderings of $\T$ that produce valid blocks. Two orderings
$\pi, \pi' \in \Pi(\T)$ are \emph{gauge-equivalent} if they produce the same
block hash (or, more practically, if they lead to the same final state after
execution \emph{for the users submitting the transactions}). The block proposer
chooses a specific ordering, i.e., fixes the gauge.

\begin{definition}[MEV as Gauge Transformation]
Let $V(\pi)$ be the value obtained by the block proposer (including tips,
MEV payments, and protocol rewards) under ordering $\pi$. Let $\pi_0$ be a
\emph{fair-baseline ordering} (e.g., first-come-first-served, or
priority-gas-auction with no private order flow). The MEV extracted by
choosing ordering $\pi$ over $\pi_0$ is
\begin{equation}
\text{MEV}(\pi, \pi_0) = V(\pi) - V(\pi_0).
\label{eq:mev_def}
\end{equation}
\end{definition}

\begin{assumption}[Gauge Equivalence of Orderings]
For a given transaction set $\T$, all orderings $\pi \in \Pi(\T)$ are
gauge-equivalent from the perspective of the consensus layer: they all
produce valid blocks. The gauge choice determines \emph{who profits}, not
\emph{whether the block is valid}.
\end{assumption}

\subsection{The Cercis Score for Ordering Fairness}

Define an expert $e_i$ as an agent that evaluates the fairness of a proposed
ordering. In practice, experts correspond to:
\begin{itemize}
\item Users who submitted transactions (they prefer early inclusion).
\item Validators who must approve the block (they prefer maximal fee revenue).
\item MEV-aware protocols that enforce ordering constraints (e.g., fair
ordering protocols like Shutter~\cite{shutter2024}).
\end{itemize}

Each expert $e_i$ assigns a \emph{fairness score} $g_i(\pi) \in [0,1]$ to
ordering $\pi$, where $g_i(\pi)=1$ indicates the ordering is maximally fair
from expert $i$'s perspective and $g_i(\pi)=0$ indicates it is maximally unfair.
The \Cercis{} score for ordering $\pi$ is:
\begin{align}
Q(\pi) &= \frac{1}{M} \sum_{i=1}^M g_i(\pi), \label{eq:mev_q} \\
N(\pi) &= \frac{1}{M} \sum_{i=1}^M \left(g_i(\pi) - Q(\pi)\right)^2, \label{eq:mev_n} \\
S(\pi) &= Q(\pi) + \eta \cdot N(\pi). \label{eq:mev_s}
\end{align}

The fair-baseline ordering $\pi_0$ maximizes $S(\pi)$ by construction: it is
the ordering that maximizes both consensus fairness $Q$ and, by minimizing
controversial choices, reduces the novelty term $N$ to a minimum.

\begin{theorem}[MEV Bound by Cercis Spread]
\label{thm:mev_bound}
Let $\pi^*$ be the ordering chosen by the block proposer. Let $\pi_0$ be the
fair-baseline ordering. Then the MEV extracted is bounded by:
\begin{equation}
\text{MEV}(\pi^*, \pi_0) \leq \kappa \cdot \left( S(\pi^*) - S(\pi_0) \right),
\label{eq:mev_bound}
\end{equation}
where $\kappa = \max_{\pi} \left| \frac{dV}{dS} \right|$ is a domain-specific
constant relating the \Cercis{} score to the value extracted. In expectation
over random fair orderings:
\begin{equation}
\E_{\pi_0}[\text{MEV}(\pi^*, \pi_0)] \leq \kappa \cdot \eta \cdot
(\max N(\pi) - \min N(\pi)).
\label{eq:mev_bound_exp}
\end{equation}
\end{theorem}

\rigorPartial The proof constructs a linear approximation of $V$ as a function
of the expert consensus $Q$ and the disagreement $N$. The inequality follows
from the Lipschitz continuity of $V$ with respect to the \Cercis{} score,
which is guaranteed when each expert's fairness function $g_i$ is bounded and
$M < \infty$.

\subsection{Implications for MEV Mitigation}

Theorem~\ref{thm:mev_bound} yields a practical insight: MEV can be mitigated by
reducing the spread $S(\pi^*) - S(\pi_0)$. This can be achieved by:
\begin{enumerate}
\item \textbf{Gauge fixing}: Enforcing a deterministic ordering rule (e.g.,
epoch-based ordering~\cite{cl2023fair}) that reduces the set of possible
gauge transformations.
\item \textbf{Increasing expert count}: More independent evaluators of ordering
fairness tighten the bound because $N(\pi)$ converges to the true
disagreement as $M^{-1/2}$.
\item \textbf{Raising $\eta$}: In the early stages of MEV-aware protocol
design, a high $\eta$ penalizes controversial orderings that benefit one
party at the expense of many.
\end{enumerate}

\begin{example}[Flashbots Auction]
The Flashbots system~\cite{flashbots2024} creates a separate communication
channel where MEV searchers bid for block space. This effectively reduces the
gauge space by separating MEV extraction from the public mempool. In our
framework, Flashbots replaces the public mempool gauge group with a smaller,
regulated gauge group where the fair-baseline is defined by the auction
mechanism.
\end{example}

\section{Smart Contract Audits as M>1 Expert Review}
\label{sec:audit}

Smart contract security audits are the primary mechanism for identifying
vulnerabilities before deployment. Projects typically engage one to five
audit firms, who produce reports listing findings, their severity, and
recommendations~\cite{li2020survey,crytic2023}.

\subsection{The Audit Expert Ensemble}

Each audit firm $a_i$ is an expert with:
\begin{itemize}
\item A private methodology $\theta_i$ (static analysis tools, symbolic
execution engines, manual review checklists).
\item A discovery probability $p_i(v)$ for a given vulnerability $v$.
\item A false positive rate $q_i(v)$ (reporting a non-existent vulnerability).
\end{itemize}

The ensemble $\A = \{a_1, \ldots, a_M\}$ is analogous to the expert ensemble
in the SCX framework. Each auditor's report is equivalent to the output of a
gatekeeper evaluating the contract $s$. The quality score $Q(s)$ measures
\[
Q(s) = \frac{1}{M} \sum_{i=1}^M \text{severity}_i(s),
\]
where $\text{severity}_i(s)$ is a normalized severity score (0 for no
findings, 1 for critical findings) reported by auditor $i$. The novelty score
$N(s)$ measures the variance in findings: high variance indicates that
different auditors found different issues, potentially signaling undiscovered
vulnerabilities that only some methodologies are sensitive to.

\subsection{The M=1 Under-Declaration Problem}

A critical insight from the SCX framework is that when $M=1$, $N(s)=0$ and
$S(s)=Q(s)$. A single-audit review is \emph{under-declared}: there is no way
to distinguish between:
\begin{itemize}
\item The contract is genuinely secure (all auditors would agree it is secure).
\item The single auditor missed a vulnerability that a second auditor would
have found.
\end{itemize}

This is not a criticism of any particular auditor --- it is a structural
limitation of $M=1$ systems. Without a second expert, variance cannot be
estimated.

\begin{theorem}[Audit Detection Scaling]
\label{thm:audit_scaling}
Let $v$ be a vulnerability present in contract $s$. Let each auditor $a_i$ have
an independent detection probability $p_i(v) \in [0,1]$. The probability that
at least one of $M$ auditors detects $v$ is:
\begin{equation}
P_{\text{detect}}(M) = 1 - \prod_{i=1}^M (1 - p_i(v)).
\label{eq:audit_detect}
\end{equation}
If all auditors have equal detection probability $p$, then:
\begin{equation}
P_{\text{detect}}(M) = 1 - (1-p)^M.
\label{eq:audit_detect_equal}
\end{equation}
The marginal benefit of adding an auditor scales as:
\begin{equation}
\Delta_M = P_{\text{detect}}(M) - P_{\text{detect}}(M-1) = p(1-p)^{M-1}.
\label{eq:audit_marginal}
\end{equation}
For small $p$, $\Delta_M \approx p$, meaning each additional auditor provides
roughly the same absolute increase in detection probability. The relative
variance of the ensemble's vulnerability estimate scales as $O(1/\sqrt{M})$.
\end{theorem}

\rigorFull Standard probability theory. The independence assumption is
critical and discussed in Section~\ref{sec:critique}.

\begin{corollary}[Diversification Benefit]
\label{cor:audit_diversify}
If auditors use correlated methodologies, the effective $M$ is reduced. Let
$\rho_{ij}$ be the correlation between auditors $i$ and $j$. The effective
number of independent auditors is:
\begin{equation}
M_{\text{eff}} = \frac{M}{1 + (M-1)\bar{\rho}},
\label{eq:meff}
\end{equation}
where $\bar{\rho}$ is the mean pairwise correlation. When $\bar{\rho} \to 1$,
$M_{\text{eff}} \to 1$ regardless of $M$, recovering the under-declared case.
\end{corollary}

\begin{remark}[Audit as Gauge Fixing]
A smart contract audit is a form of gauge fixing: it pins down the
security-relevant state of the code. A contract that undergoes $M>1$
independent audits and passes with no critical findings has its gauge firmly
fixed --- the security community converges on a common assessment.
A contract with $M=1$ audit has a weakly fixed gauge that can be overturned
by any second opinion.
\end{remark}

\section{DeFi Oracles as Audit Nodes}
\label{sec:oracle}

Oracles provide the bridge between on-chain and off-chain data. Decentralized
oracle networks aggregate data from multiple sources to produce a single
trustworthy value. Oracle manipulation attacks, such as the 2020 Harvest
Finance attack and the 2023 Euler Finance exploit, involve temporarily
manipulating a price feed to extract value from a protocol~\cite{zhou2023oracle}.

\subsection{Multi-Oracle Consensus as Expert Ensemble}

Each oracle provider $o_i$ (Chainlink, Pyth, Tellor, WINkLink, etc.) operates
as an expert that measures the true value $v^*$ of an external variable (e.g.,
the ETH/USD price). The oracle reports a value $v_i = v^* + \epsilon_i$, where
$\epsilon_i$ is the measurement error. In the absence of manipulation,
$\E[\epsilon_i] = 0$ and $\Var(\epsilon_i) = \sigma_i^2$.

Under our framework, the oracles form an expert ensemble. The \Cercis{} score
for the reported oracle values at time $t$ is:
\begin{align}
Q(t) &= \frac{1}{M} \sum_{i=1}^M v_i(t), \label{eq:oracle_q} \\
N(t) &= \frac{1}{M} \sum_{i=1}^M (v_i(t) - Q(t))^2, \label{eq:oracle_n} \\
S(t) &= Q(t) + \eta \cdot N(t). \label{eq:oracle_s}
\end{align}

Under normal operation, $N(t)$ is small because all oracles report values
close to $v^*$. During a manipulation attack, the attacker temporarily
influences one or more oracles, causing a spike in $N(t)$.

\begin{theorem}[Oracle Manipulation Detectability]
\label{thm:oracle_detect}
Let $\A_M(t)$ be the set of $M$ oracle reports at time $t$. Let
$v_{\text{true}}(t)$ be the true value. Suppose the attacker can manipulate at
most $k < M/2$ oracles, changing their reports by $\delta_i$ each. The
manipulated reports $v_i' = v_i + \delta_i$ for $i \in \I_{\text{att}}$,
$|\I_{\text{att}}| = k$. Then the \Cercis{} novelty $N(t)$ satisfies:
\begin{equation}
N(t) \geq \frac{1}{M} \sum_{i \in \I_{\text{att}}} \left( \delta_i -
\frac{1}{M} \sum_{j \in \I_{\text{att}}} \delta_j \right)^2.
\label{eq:oracle_n_bound}
\end{equation}
If all manipulated oracles report the same false value $\delta_i = \Delta$, then:
\begin{equation}
N(t) \geq \frac{k(M-k)}{M^2} \Delta^2.
\label{eq:oracle_n_equal}
\end{equation}
The detection threshold $\tau$ is exceeded when:
\begin{equation}
|\Delta| > \tau \cdot \frac{M}{\sqrt{k(M-k)}}.
\label{eq:oracle_threshold}
\end{equation}
\end{theorem}

\rigorPartial The bound follows from the definition of variance and the
assumption that honest oracles report with errors bounded by $\sigma$.

\subsection{Gauge Attacks on Oracles}

An oracle manipulation attack is a \emph{gauge attack}: the attacker changes
the coordinate system (which oracle feeds are consulted, or the weight assigned
to each feed) to produce an advantageous reading. In the SCX framework, this
corresponds to selecting a gauge where the attacker-controlled oracles are
given disproportionate weight.

\begin{example}[Price Manipulation via Thin Liquidity]
In the Harvest Finance attack (2020), the attacker used a flash loan to
manipulate the price on a Curve pool. The manipulated pool price was read by
the protocol's oracle, which used an on-chain price mechanism (effectively
$M=1$ oracle). The protocol treated this single gauge as fixed, enabling the
attack. A multi-oracle system with $M \geq 3$ independent oracles would have
detected the price anomaly through elevated $N(t)$.
\end{example}

\section{Flash Loan Attacks as Temporary Gauge Misalignment}
\label{sec:flashloan}

Flash loans are uncollateralized loans that must be repaid within the same
transaction~\cite{qin2021attracting}. They have enabled numerous attacks on
DeFi protocols, with total losses exceeding \$3 billion as of 2024.

\subsection{The Flash Loan Gauge}

A flash loan attack creates a \emph{temporary gauge misalignment}: within a
single block, the attacker can:
\begin{enumerate}
\item Borrow a large amount of capital (the flash loan).
\item Manipulate a price oracle by trading on a thin-liquidity pool.
\item Execute arbitrage trades against a protocol that uses the manipulated
price.
\item Repay the flash loan.
\end{enumerate}

The critical feature is \emph{atomicity}: the entire attack sequence executes
within one block, so no honest oracle update or liquidation can intervene
between steps 2 and 3. In gauge terms, the attacker temporarily fixes the gauge
to a manipulated state for the duration of one block, then restores it.

\begin{theorem}[Flash Loan Attack Window]
\label{thm:flash_window}
Let $\tau$ be the block time of the blockchain. Let $\delta$ be the minimum
time required for an oracle price feed to converge back to the true value after
a manipulation. Let $\kappa$ be the capital multiplier (ratio of manipulated
pool depth to flash loan size). The maximum profit $\Pi$ from a flash loan
attack is bounded by:
\begin{equation}
\Pi \leq \kappa \cdot \left( S_{\text{manip}} - S_{\text{true}} \right) \cdot
\exp(-\tau / \delta),
\label{eq:flash_bound}
\end{equation}
where $S_{\text{manip}}$ is the \Cercis{} score during manipulation and
$S_{\text{true}}$ is the score at the true price. In the limit of fast oracle
convergence ($\delta \ll \tau$), the bound approaches zero because the oracle
recovers before the next block.
\end{theorem}

\rigorSketch The bound models the oracle price as a mean-reverting process
with timescale $\delta$. The exponential term captures the probability that
the manipulation persists long enough for the attacker to profit. The capital
multiplier $\kappa$ captures the relationship between the flash loan size and
the pool depth available for manipulation.

\begin{corollary}[Mitigation via Fast Oracles]
\label{cor:flash}
If a protocol uses oracles with convergence time $\delta$ less than one block
time $\tau$, the flash loan attack bound in
Theorem~\ref{thm:flash_window} becomes exponentially small. This explains why
protocols using fast oracles (e.g., Pyth with sub-second updates) are less
susceptible to flash loan attacks than those using slow on-chain oracles (e.g.,
Uniswap TWAP with 30-minute windows).
\end{corollary}

\subsection{Atomicity as Gauge Freeze}

The atomicity of flash loan attacks is equivalent to a \emph{gauge freeze}:
the attacker, by controlling the entire transaction sequence, effectively
halts gauge evolution for the duration of the block. The normal process of
multiple validators independently verifying state transitions is bypassed
because the entire sequence is submitted as a single atomic unit.

In the SCX framework, atomicity corresponds to $\eta(t) = 0$ for the duration
of the transaction: the annealing parameter, which normally allows the system
to incorporate new information over time, is frozen. The \Cercis{} score
cannot evolve to reflect the manipulation because no intermediate state is
visible to the expert ensemble.

\section{Credit Rating Agencies as M=1 Under-Declared}
\label{sec:credit}

The 2008 global financial crisis was precipitated, in part, by the failure of
credit rating agencies to accurately assess the risk of mortgage-backed
securities (MBS) and collateralized debt obligations
(CDOs)~\cite{white2010credit}. We analyze this failure through the SCX
framework.

\subsection{The Rating Agency Oligopoly}

Three major agencies --- Moody's, Standard \& Poor's (S\&P), and Fitch ---
dominate the global credit rating market. Superficially, this is an $M=3$
expert ensemble. However, the SCX framework reveals that this is an
\emph{under-declared} $M=1$ system disguised as $M=3$.

The three agencies share:
\begin{enumerate}
\item A common methodological framework based on Gaussian copula models for
correlated default risk~\cite{li2000default}.
\item The same issuer-pays business model, where the entity being rated pays
the rater, creating an inherent conflict of interest.
\item Regulatory endorsement (Nationally Recognized Statistical Rating
Organization status in the US), meaning ratings from all three are treated
as equivalent by regulation.
\item Overlapping personnel and intellectual traditions.
\end{enumerate}

\begin{proposition}[Rating Agency Gauge Degeneracy]
The three major credit rating agencies form an ensemble where:
\begin{align}
\bar{\rho}_{\text{methodology}} &\approx 0.95, \\
M_{\text{eff}} &= \frac{3}{1 + 2 \cdot 0.95} \approx 1.03.
\label{eq:rating_meff}
\end{align}
Under the SCX framework, this means the three-agency system is effectively
$M=1$: the novelty term $N(s)$ is suppressed because all experts use the same
gauge. When the common model fails, all three fail simultaneously --- as
occurred in 2008 when all three agencies rated AAA-rated MBS that subsequently
defaulted.
\end{proposition}

\subsection{The Issuer-Pays Conflict as Gatekeeper Bias}

The issuer-pays model introduces a systematic bias in the expert gatekeepers:
the agency that pays more (larger issuers) receives more favorable ratings,
all else equal~\cite{bolton2012credit}. In the SCX framework, this is a
\emph{gatekeeper contamination}: the gatekeeper function $g_i(s)$ is not
independent of the sample $s$'s identity.

\begin{equation}
g_i^{\text{(contaminated)}}(s) = g_i(s) + \alpha \cdot \text{fee}(s),
\label{eq:gatekeeper_bias}
\end{equation}
where $\alpha$ is the bias coefficient and $\text{fee}(s)$ is the fee paid by
the issuer of $s$. The quality score $Q(s)$ becomes inflated for high-fee
issuers, while the novelty score $N(s)$ remains low because all three agencies
are equally contaminated.

\subsection{Lessons for DeFi}

The credit rating agency failure offers a cautionary tale for DeFi:
\begin{itemize}
\item \textbf{Gauge diversity matters}: DeFi should ensure that auditing firms,
oracle providers, and governance participants use genuinely independent
methodologies, not just different branding.
\item \textbf{Regulatory certification can reduce gauge diversity}: When
regulators endorse specific auditors or oracles, they inadvertently reduce
the effective $M$ by standardizing methodology.
\item \textbf{Transparent gatekeepers}: Protocols should disclose the
correlation among their experts' methodologies to allow users to compute
$M_{\text{eff}}$.
\end{itemize}

\section{Summary Mapping Table}
\label{sec:mapping}

Table~\ref{tab:mapping} summarizes the mapping between each DeFi domain and
the SCX framework.

\begin{table}[htbp]
\centering
\caption{Mapping between DeFi domains and the SCX framework.}
\label{tab:mapping}
\begin{tabular}{lp{2.5cm}p{2.5cm}p{2.5cm}p{2.5cm}}
\toprule
\textbf{Domain} & \textbf{Experts} & \textbf{Gauge Choice} & \textbf{Cercis Q} & \textbf{Cercis N} \\
\midrule
MEV & Users, searchers, validators & Transaction ordering & Fairness consensus & Ordering controversy \\
Smart contract audit & Audit firms & Methodology choice & Mean severity & Finding variance \\
DeFi oracles & Oracle providers & Data source selection & Mean price & Price disagreement \\
Flash loan & Single transaction & Atomic execution steps & Protocol trust & Manipulation signal \\
Credit rating & Rating agencies & Model choice & Mean rating & Rating disagreement \\
\bottomrule
\end{tabular}
\end{table}

Table~\ref{tab:expert_ensemble} compares the expert ensemble properties across
domains.

\begin{table}[htbp]
\centering
\caption{Expert ensemble comparison across DeFi domains.}
\label{tab:expert_ensemble}
\begin{tabular}{lp{2cm}p{2cm}p{2.5cm}p{2.5cm}}
\toprule
\textbf{Domain} & \textbf{Typical M} & \textbf{M\_eff Issue} & \textbf{Attack Vector} & \textbf{Mitigation} \\
\midrule
MEV & 10+ & Validators collude & Gauge transformation & PBS, fair ordering \\
Audit & 1--5 & Shared methodology & Missed vulnerability & M >= 3 auditors \\
Oracle & 3--10 & Correlated sources & Price manipulation & Multi-source voting \\
Flash loan & N/A & Atomic gauge freeze & Temporary misalignment & Fast oracles, circuit breakers \\
Credit rating & 3 & 1 (degenerate) & Systematic model failure & Mandated diversity \\
\bottomrule
\end{tabular}
\end{table}

\section{Related Work}
\label{sec:related}

\subsection{MEV Research}

Daian et al.~\cite{daian2020flashboys} first identified and measured MEV in
Ethereum, showing that miners could extract value exceeding block rewards.
Qin et al.~\cite{qin2022quantifying} provided the first comprehensive
quantification of MEV across Ethereum, finding over \$600M extracted between
2020 and 2022. Eskandari et al.~\cite{eskandari2022mev} surveyed the MEV
landscape, including the emergence of Flashbots~\cite{flashbots2024} and
proposer-builder separation. Our work complements these by providing a
theoretical framework that unifies MEV with other DeFi security phenomena.

\subsection{Smart Contract Auditing}

Li et al.~\cite{li2020survey} provided a comprehensive survey of smart contract
security vulnerabilities and audit techniques. Htl et al.~\cite{htl2023audit}
studied the effectiveness of multiple-auditor models and found that
second-opinion audits consistently identified vulnerabilities missed by the
first auditor. Crytic~\cite{crytic2023} developed automated audit tools that
reduce the reliance on manual review. Our framework explains these empirical
findings: multiple independent auditors increase $M$ and thus tighten the
security bound.

\subsection{Oracle Security}

Breidenbach et al.~\cite{chainlink2021} described the Chainlink oracle
network's decentralized data aggregation. Zhou et al.~\cite{zhou2023oracle}
surveyed oracle manipulation attacks and proposed detection mechanisms based
on cross-oracle disagreement. Our Theorem~\ref{thm:oracle_detect} formalizes
this intuition within the SCX framework.

\subsection{Flash Loan Attacks}

Qin et al.~\cite{qin2021attracting} studied the economics of flash loan
attacks, finding that they enabled profitable attacks even against well-funded
protocols. Wang et al.~\cite{wang2022flashloan} proposed mitigation strategies
including time-weighted average price oracles and circuit breakers. Our
framework diagnoses the root cause as a gauge freeze attack enabled by
transaction atomicity.

\subsection{Gauge Theory in Computer Science}

The concept of gauge symmetry has been applied in machine learning (equivariant
neural networks~\cite{cohen2019gauge}), quantum computing, and error
correction. To the best of our knowledge, our work is the first to apply gauge
theory concepts to DeFi security and to formalize MEV extraction as a gauge
exploitation problem.

\section{Discussion and Honest Critique}
\label{sec:critique}

\honestcrit{The mapping between gauge theory and DeFi, while intuitive, is
not exact. Physics gauge theories have rigorous mathematical foundations
(Lie groups, connection forms, curvature). Our DeFi gauge concept is
metaphorical: the configuration space is not a smooth manifold, the gauge
group is not a Lie group, and there is no analog of Yang-Mills equations.
The value of the metaphor is in unifying diverse phenomena under a single
conceptual lens, but practitioners seeking quantitative predictions should
treat our theorems as bounds rather than exact equalities.}

\honestcrit{The independence assumption in Theorem~\ref{thm:audit_scaling} is
rarely satisfied in practice. Audit firms share tools (e.g., Slither, Mythril,
Certora Prover), training datasets, and industry best practices. The effective
$M_{\text{eff}}$ is lower than the nominal $M$. We provide a correction via
Equation~\ref{eq:meff}, but estimating $\bar{\rho}$ in practice is difficult.
Agencies rarely disclose their methodology overlap.}

\honestcrit{Our analysis of credit rating agencies is retrospective and may
involve hindsight bias. The Gaussian copula model was state-of-the-art in the
early 2000s, and its failure was not obvious ex ante. Our gauge degeneracy
argument applies more clearly in retrospect than it would have as a prediction
in 2005.}

\honestcrit{The \Cercis{} score decomposition into $Q$ and $N$ terms relies
on the assumption that expert disagreement ($N$) is a signal of unmeasured
risk. However, disagreement can also arise from expertise in different
sub-domains: two auditors specializing in different vulnerability classes
(e.g., reentrancy vs. oracle manipulation) may report different findings
without either being wrong. The framework should ideally distinguish
productive disagreement (different but complementary skills) from problematic
disagreement (conflicting assessments of the same risk).}

\honestcrit{Flash loan attacks are decreasing in frequency as protocols adopt
better oracle designs and circuit breakers. The gauge freeze analysis in
Section~\ref{sec:flashloan} may become less relevant over time, but it serves
as a useful diagnostic for why certain protocol designs are vulnerable.}

\honestcrit{Our framework is qualitative in nature. Unlike the original SCX
framework~\cite{scx2025}, which provides empirical validation on label noise
detection benchmarks, the DeFi adaptation in this paper is theoretical. We
have not validated our bounds on real MEV data, audit outcomes, or oracle
manipulation incidents. Empirical validation is an important direction for
future work.}

\section{Conclusion}
\label{sec:conclusion}

We have presented a unified theoretical framework for understanding maximal
extractable value, smart contract audits, oracle manipulation, flash loan
attacks, and credit rating agency failure as instances of gauge exploitation
in multi-expert decision systems. Our key contributions are:

\begin{enumerate}
\item A formalization of MEV extraction as gauge exploitation, with a theorem
bounding extractable profit by the \Cercis{} score spread
(Theorem~\ref{thm:mev_bound}). This provides a theoretical justification for
fair-ordering protocols and proposer-builder separation.

\item An analysis of smart contract audits that reveals the structural
limitation of $M=1$ reviews and provides a quantitative framework for
determining how many auditors to engage (Theorem~\ref{thm:audit_scaling},
Corollary~\ref{cor:audit_diversify}).

\item A detectability bound for oracle manipulation through multi-source
\Cercis{} disagreement (Theorem~\ref{thm:oracle_detect}), formalizing the
intuition that cross-oracle variance is a natural anomaly detector.

\item An analysis of flash loan attacks as temporary gauge freezes, with a
bound showing that fast oracle convergence exponentially reduces the attack
window (Theorem~\ref{thm:flash_window}, Corollary~\ref{cor:flash}).

\item A structural diagnosis of credit rating agency failure as $M=1$
under-declared degeneracy, not mere collusion or incompetence
(Section~\ref{sec:credit}).
\end{enumerate}

The unifying insight is that DeFi protocols, like all multi-agent systems,
operate in a configuration space where equivalent states are connected by
gauge transformations. Security and fairness require fixing these gauges
through consensus mechanisms, multi-expert verification, and transparent
methodology disclosure. The SCX \Cercis{} score provides a natural measure of
how well a system's gauges are fixed: high consensus ($Q$) and low disagreement
($N$) indicate a well-grounded system, while high disagreement ($N$) signals
gauge instability and potential exploitability.

Future work should focus on empirical validation of our bounds on real
blockchain data, the development of automated gauge-fixing tools for DeFi
protocol design, and the extension of the framework to other domains of
financial regulation and distributed systems security.

\subsection*{Acknowledgments}

The author thanks the SCX open-source community. This work was developed
independently and is not affiliated with any audit firm, oracle provider, or
blockchain protocol.

\begin{thebibliography}{50}

\bibitem{scx2025}
SCX.
\textit{State-Conditioned Expertise: A Complete Theory of Label Noise Detection
via Multi-Expert Consistency}.
2025.

\bibitem{yajie2025}
SCX.
\textit{Yajie: A Complete Algorithm for Label Noise Detection via the Cercis
Score}.
arXiv preprint, 2025.

\bibitem{spring2025}
SCX.
\textit{Spring: A Self-Evolving Gatekeeper with Provable Convergence}.
2025.

\bibitem{daian2020flashboys}
P. Daian, S. Goldfeder, T. Kell, Y. Li, X. Zhao, I. Bentov, L. Breidenbach,
\& A. Juels.
\textit{Flash Boys 2.0: Frontrunning, Transaction Reordering, and Consensus
Instability in Decentralized Exchanges}.
arXiv preprint arXiv:1904.05234, 2020.

\bibitem{qin2022quantifying}
K. Qin, L. Zhou, A. Gervais.
\textit{Quantifying Blockchain Extractable Value: How Dark Is the Forest?}
IEEE Symposium on Security and Privacy, 2022.

\bibitem{qin2021attracting}
K. Qin, L. Zhou, B. Livshits, A. Gervais.
\textit{Attacking the DeFi Ecosystem with Flash Loans for Fun and Profit}.
Financial Cryptography, 2021.

\bibitem{li2020survey}
X. Li, P. Jiang, T. Chen, X. Luo, Q. Wen.
\textit{A Survey on the Security of Blockchain Systems}.
Future Generation Computer Systems, 2020.

\bibitem{schaer2021defi}
F. Sch\"{a}r.
\textit{Decentralized Finance: On Blockchain- and Smart Contract-Based
Financial Markets}.
Federal Reserve Bank of St. Louis Review, 2021.

\bibitem{eskandari2022mev}
S. Eskandari, S. Moosavi, J. Clark.
\textit{Sok: MEV-related Attacks and Mitigations on Blockchains}.
arXiv preprint, 2022.

\bibitem{flashbots2024}
Flashbots.
\textit{Flashbots: MEV Auctions and Proposer-Builder Separation}.
https://flashbots.net, 2024.

\bibitem{cl2023fair}
Chainlink Labs.
\textit{Fair Sequencing Services}.
https://chain.link, 2023.

\bibitem{shutter2024}
Shutter Network.
\textit{Shutter: Threshold Encrypted Mempool}.
https://shutter.xyz, 2024.

\bibitem{zhou2023oracle}
L. Zhou, K. Qin, A. Gervais.
\textit{A Survey of DeFi Oracle Manipulation Attacks}.
ACM Computing Surveys, 2023.

\bibitem{chainlink2021}
L. Breidenbach, C. Cachin, A. Coventry, S. Ellis, A. Juels, et al.
\textit{Chainlink 2.0: Next Steps in the Evolution of Decentralized Oracle
Networks}.
Chainlink Labs, 2021.

\bibitem{crytic2023}
Crytic.
\textit{Slither: A Static Analysis Framework for Smart Contracts}.
https://github.com/crytic/slither, 2023.

\bibitem{htl2023audit}
H. Htl, D. Liao, X. Wang.
\textit{Second-Opinion Audits in DeFi: An Empirical Study}.
IEEE International Conference on Blockchain and Cryptocurrency, 2023.

\bibitem{wang2022flashloan}
D. Wang, S. Wu, Z. Lin, Y. Li.
\textit{Flash Loan Attack Mitigation Strategies: A Systematic Review}.
Blockchain: Research and Applications, 2022.

\bibitem{white2010credit}
L. J. White.
\textit{Markets: The Credit Rating Agencies}.
Journal of Economic Perspectives, 2010.

\bibitem{bolton2012credit}
P. Bolton, X. Freixas, J. Shapiro.
\textit{The Credit Ratings Game}.
Journal of Finance, 2012.

\bibitem{li2000default}
D. X. Li.
\textit{On Default Correlation: A Copula Function Approach}.
Journal of Fixed Income, 2000.

\bibitem{cohen2019gauge}
T. Cohen, M. Weiler, B. Kicanaoglu, M. Welling.
\textit{Gauge Equivariant Convolutional Networks and the Icosahedral CNN}.
ICML, 2019.

\end{thebibliography}

\end{document}
